\documentclass[11pt,a4paper]{article}

% ─────────── Packages ───────────
\usepackage[utf8]{inputenc}
\usepackage[spanish]{babel}
\usepackage[T1]{fontenc}
\usepackage{ebgaramond}            % Serif elegante
\usepackage[margin=2.7cm]{geometry}
\usepackage{microtype}
\usepackage{graphicx}
\usepackage{booktabs}
\usepackage{multirow}
\usepackage{amsmath, amssymb}
\usepackage{siunitx}
\usepackage[hidelinks]{hyperref}
\usepackage[round, sort&compress]{natbib}
\usepackage{authblk}
\usepackage{titlesec}
\usepackage{xcolor}

\definecolor{accent}{HTML}{7A1F1F}
\definecolor{ink}{HTML}{1A1612}
\hypersetup{citecolor=accent, urlcolor=accent, linkcolor=ink}

\titleformat*{\section}{\Large\itshape\color{ink}}
\titleformat*{\subsection}{\large\itshape\color{ink}}

% ─────────── Front matter ───────────
\title{\textbf{El millón invisible}\\[0.3em]
\large\itshape Divergencia rama--ocupación en el empleo tecnológico español\\
y la performatividad de la estadística pública (EPA T1 2026)}

\author[1]{José Fernández Tamames\thanks{
\href{mailto:jtamames@unie.es}{jtamames@unie.es};
ORCID: \href{https://orcid.org/0009-0007-8851-9833}{0009-0007-8851-9833}}}
\affil[1]{UNIE Universidad, Madrid · Cátedra de Filosofía y Tecnología}

\date{29 de abril de 2026 \\ \textit{Borrador para \textit{Ratio \& Machina}, vol. I, n.\textsuperscript{o} 1}}

\begin{document}
\maketitle

% ─────────── Abstract ───────────
\begin{abstract}
\noindent
La medición pública del empleo tecnológico en España convive con dos definiciones
operativas que arrojan magnitudes radicalmente distintas: una basada en la rama
de actividad (Sección J de la CNAE-2009) y otra en la ocupación (códigos CNO-2011
de profesionales TIC). En 2024, la primera contabiliza unos 644.000 ocupados; la
segunda supera el millón. La diferencia ---el ``millón invisible''--- corresponde
a profesionales tecnológicos embebidos en sectores no tecnológicos. Sobre la
serie reconstruida 2019Q3--2026Q1, el test de Mann-Kendall confirma tendencia
creciente significativa del agregado ($Z = 5{,}74$, $p < 0{,}001$) con
recomposición interna heterogénea: caída estructural de Telecomunicaciones,
crecimiento sostenido de Actividades informáticas. La descomposición STL
demuestra estacionalidad cuasi-nula del sector, anulando el argumento del ruido
estacional para explicar contracciones trimestrales. Más allá de los hallazgos
empíricos, el trabajo propone una lectura filosófica: la transición CNAE-2009
$\rightarrow$ CNAE-2025 (Reglamento Delegado UE 2023/137), que escinde la
Sección J en dos secciones nuevas, materializa la \emph{performatividad} del
dato administrativo. Las ramas más expuestas a la automatización
---almacenamiento, reparación de maquinaria, actividades técnicas auxiliares---
cambian de hogar estadístico precisamente cuando más relevante sería medir su
evolución con consistencia. La estadística pública es ya un actor del fenómeno
que mide.

\medskip
\noindent\textbf{Palabras clave:} empleo tecnológico, EPA, CNAE, performatividad
del dato, filosofía de la tecnología, causa segunda, automatización.

\medskip
\noindent\textbf{Códigos JEL:} J21, J24, O33, C82.
\end{abstract}

% ═══════════════════════════════════════════════════════════════════════════
\section{Introducción}

La pregunta sobre cuántas personas trabajan en el ámbito tecnológico en España
parece resoluble por consulta directa a la fuente oficial: la Encuesta de
Población Activa (EPA) del Instituto Nacional de Estadística publica
trimestralmente el empleo desagregado por rama de actividad económica. Sin
embargo, dos cifras concurren legítimamente, sin que ninguna sea errónea, para
responder a esa pregunta. La primera ---rama de actividad--- contabiliza
ocupados en empresas cuya actividad principal está clasificada como
``Información y comunicaciones'' (Sección J de la CNAE-2009). La segunda
---ocupación--- contabiliza personas cuyo puesto desempeñado se codifica
como ``profesional TIC'' en la CNO-2011, con independencia del sector económico
de la empresa empleadora. Las dos tienen respaldo metodológico en estadísticas
oficiales (EPA, ONTSI) y agregaciones europeas (Eurostat).

La discrepancia entre ambas, lejos de ser un artefacto técnico, expresa una
transformación estructural del trabajo: la difusión del trabajo tecnológico
\emph{fuera} del perímetro de las empresas tecnológicas. Esta nota empírica
documenta el fenómeno para España (2019--2026), formaliza estadísticamente su
crecimiento, y propone una lectura filosófica de sus implicaciones para la
estadística pública en el contexto de la transición CNAE-2009 $\rightarrow$
CNAE-2025.

% ═══════════════════════════════════════════════════════════════════════════
\section{Datos y método}

\subsection{Fuentes}

Se utilizan tres fuentes complementarias:
\begin{itemize}
\item \textbf{EPA-INE}, notas de prensa trimestrales 2019Q3--2026Q1 y series
  agregadas de la Sección J reconstruidas a partir de la edición Randstad
  Research \citep{randstad2026}, que desagrega Sección J en sus tres divisiones
  (61 Telecomunicaciones, 62 Programación/consultoría, 63 Servicios de
  información).
\item \textbf{ONTSI 2025} \citep{ontsi2025}, informe anual sobre empleo
  tecnológico que reporta especialistas TIC por ocupación (CNO 25/35) con serie
  2018--2024.
\item Reglamento Delegado UE 2023/137 \citep{ue2023137} sobre la transición
  NACE Rev. 2 $\rightarrow$ NACE Rev. 2.1 / CNAE-2025 \citep{rd102025}.
\end{itemize}

\subsection{Métodos estadísticos}

\paragraph{Test de tendencia.} Se aplica el test de Mann-Kendall modificado por
\citet{hamed1998} para corregir la autocorrelación de las series trimestrales
de empleo. La pendiente de Sen estima la magnitud anual de la tendencia.

\paragraph{Descomposición.} La componente tendencial, estacional y residual se
estiman mediante STL (Seasonal-Trend decomposition using LOESS) con período
$s = 4$ y ajuste robusto. Se reportan ratios $\text{Var}(C)/\text{Var}(Y)$ para
$C \in \{T, S, R\}$.

\paragraph{Comparación entre subsectores.} Test no paramétrico de Friedman
sobre las distribuciones de variaciones interanuales por subsector ($n=23$
trimestres). Análisis post-hoc por pares con Wilcoxon de rangos signados y
corrección Bonferroni ($\alpha = 0{,}05/3 = 0{,}0167$).

% ═══════════════════════════════════════════════════════════════════════════
\section{Resultados}

\subsection{La divergencia rama/ocupación}

La Tabla~\ref{tab:divergencia} muestra la evolución del gap entre ambas
métricas. En 2019, los especialistas TIC superaban en 192.600 personas al
empleo en Sección J ($+40{,}8\%$ sobre la rama); en 2024, el gap asciende a
355.800 ($+55{,}2\%$). El test de Mann-Kendall sobre la serie del gap absoluto
arroja $Z=2{,}73$, $p=0{,}006$, con pendiente de Sen $\approx +28$~mil
personas/año.

\begin{table}[h]
\centering
\caption{Evolución de la divergencia rama--ocupación, 2019--2024
\textit{(miles de personas)}}
\label{tab:divergencia}
\begin{tabular}{lrrrr}
\toprule
\textbf{Año} & \textbf{Rama (Sección J)} & \textbf{Ocupación (CNO TIC)} & \textbf{Gap} & \textbf{Gap \%} \\
\midrule
2019 & 472,4 & 665 & 192,6 & 40,8 \\
2020 & 485,4 & 727 & 241,6 & 49,8 \\
2021 & 516,6 & 775 & 258,5 & 50,0 \\
2022 & 555,8 & 820 & 264,2 & 47,5 \\
2023 & 589,5 & 915 & 325,5 & 55,2 \\
\textbf{2024} & \textbf{644,2} & \textbf{1.000} & \textbf{355,8} & \textbf{55,2} \\
\bottomrule
\end{tabular}
\end{table}

\subsection{Tendencia y estacionalidad}

La Tabla~\ref{tab:mk} resume los resultados del test de Mann-Kendall por
subsector. Actividades informáticas exhibe tendencia creciente significativa
($Z=2{,}83$); el agregado de Sección J, fuertemente significativa ($Z=5{,}74$).
Telecomunicaciones muestra tendencia decreciente no significativa pero
consistente. La descomposición STL (Tabla~\ref{tab:stl}) revela que la
componente estacional explica menos del 1\% de la varianza en todas las series,
muy por debajo del rango habitual en sectores estacionales como hostelería
($> 50\%$).

\begin{table}[h]
\centering
\caption{Test de Mann-Kendall modificado · series 2019Q3--2026Q1}
\label{tab:mk}
\begin{tabular}{lrrrl}
\toprule
\textbf{Serie} & $\boldsymbol{Z}$ & $\boldsymbol{p}$ & \textbf{Sen/año} & \textbf{Tendencia} \\
\midrule
Actividades informáticas & 2,83 & 0,005 & +26,2 & creciente** \\
Telecomunicaciones       & --1,55 & 0,122 & --2,4 & no signif. \\
Servicios información    & --0,50 & 0,619 & --0,8 & ambigua \\
\textbf{Total Sección J} & \textbf{5,74} & $\boldsymbol{<0{,}001}$ & \textbf{+34,4} & \textbf{creciente***} \\
\bottomrule
\multicolumn{5}{l}{\footnotesize\textit{Significatividad}: *** $p<0{,}001$; ** $p<0{,}01$; * $p<0{,}05$.}
\end{tabular}
\end{table}

\begin{table}[h]
\centering
\caption{Descomposición STL · varianza explicada por componente}
\label{tab:stl}
\begin{tabular}{lrrr}
\toprule
\textbf{Serie} & $\text{Var}(T)/\text{Var}(Y)$ & $\text{Var}(S)/\text{Var}(Y)$ & $\text{Var}(R)/\text{Var}(Y)$ \\
\midrule
Actividades informáticas & 0,972 & 0,000 & 0,003 \\
Telecomunicaciones       & 0,927 & 0,006 & 0,032 \\
Total Sección J          & 0,912 & 0,000 & 0,020 \\
\bottomrule
\end{tabular}
\end{table}

\subsection{Heterogeneidad inter-subsectorial}

El test de Friedman sobre las variaciones interanuales rechaza la hipótesis nula
de equivalencia distribucional ($\chi^2 = 35{,}2$, $p < 0{,}001$). El análisis
post-hoc indica que Actividades informáticas se distingue significativamente
tanto de Telecomunicaciones como de Servicios de información ($p < 0{,}0001$ en
ambos casos), mientras que estas dos últimas no son distinguibles entre sí
($p = 0{,}58$). La cifra agregada del sector ($-2{,}2\%$ interanual en T3 2025
\citep{randstad2026}) encubre, por tanto, una recomposición interna asimétrica.

% ═══════════════════════════════════════════════════════════════════════════
\section{Discusión}

\subsection{Implicaciones empíricas}

Los hallazgos sugieren tres tesis sobre el empleo tecnológico español:

\textbf{(1)} El \emph{sector tech} como categoría operativa pierde poder
descriptivo. El crecimiento del empleo tecnológico ocurre cada vez más fuera
del perímetro CNAE de las empresas tecnológicas. La proporción del talento TIC
embebido en sectores tradicionales pasó del 41\% (2019) al 55\% (2024). La
política industrial enfocada en ``el sector TIC'' está optimizando para una
fracción decreciente del fenómeno.

\textbf{(2)} La caída de Telecomunicaciones no es coyuntural sino estructural.
Refleja la consumación del modelo de externalización de las grandes operadoras
hacia integradores y consultoras informáticas. La estadística agregada lo
percibe como contracción donde lo que ocurre es transferencia interna de
trabajadores entre divisiones (61 $\rightarrow$ 62).

\textbf{(3)} La concentración geográfica Madrid--Cataluña ($\sim 60\%$ del
empleo TIC, índice HHI $\approx 1.684$) constituye un factor de riesgo
sistémico para la transformación digital nacional, en contradicción con el
discurso oficial de la ``España conectada''.

\subsection{La performatividad de la categoría administrativa}

\noindent
La transición CNAE-2009 $\rightarrow$ CNAE-2025, en vigor desde marzo de 2025
y con doble codificación durante 2026, escinde la antigua Sección J en dos
secciones nuevas: J (edición, radiodifusión, contenidos) y K (telecomunicaciones,
programación, consultoría TI). El sector tecnológico deja de localizarse con
una sola letra. Más críticamente, las ramas más expuestas a la automatización
---reparación de maquinaria, almacenamiento, actividades técnicas auxiliares---
se redistribuyen entre Industria, Construcción y Servicios de manera que las
series longitudinales pierden continuidad.

Esta no es una limitación técnica de la EPA, sino la expresión, en el plano de
la estadística pública, de un fenómeno conceptualmente más profundo. En la
terminología filosófica que el autor ha desarrollado en otra parte
\citep{tamames2024causa}, el dato administrativo opera como \emph{causa segunda}:
no genera autónomamente el fenómeno medido, pero al codificarlo lo configura,
y al recodificarlo lo reescribe. La transición CNAE no responde a un cambio
del mundo: responde a una decisión administrativa europea (Reglamento Delegado
UE 2023/137) que, sin embargo, modifica el objeto que la propia administración
mide. La performatividad, lejos de ser una propiedad excepcional del lenguaje,
resulta ser una característica estructural de toda categorización oficial.

La consecuencia operativa es que toda comparación interanual de empleo
sectorial T1 2027 vs.\ T1 2026 que no recurra a la matriz de transición
publicada por el INE producirá artefactos espurios. La comunidad estadística
y el discurso público disponen, durante 2026, de un breve período crítico de
doble codificación para construir series de enlace; agotado este margen, la
serie histórica del empleo tecnológico español requerirá, para fines
comparativos, una operación de retraducción que necesariamente introducirá
discrecionalidad metodológica.

% ═══════════════════════════════════════════════════════════════════════════
\section{Conclusiones y agenda de investigación}

Este trabajo ha cuantificado, con métodos estadísticos no paramétricos
robustos, una divergencia creciente entre dos definiciones legítimas del
empleo tecnológico en España: una crece a +28~mil personas/año más rápido que
la otra. El fenómeno no es una rareza de la contabilidad pública sino el
síntoma estadístico de una transformación material: el trabajo tecnológico se
ha convertido en una función distribuida en toda la economía, no en una rama
sectorial localizable.

Las extensiones naturales del análisis incluyen: \emph{(a)} replicación con
microdatos EPA brutos a partir del T1 2026 (con doble codificación), una vez
disponibles; \emph{(b)} análisis comparativo a nivel europeo con NACE Rev.~2.1
como referencia; \emph{(c)} cuantificación del ``drift'' categórico inducido
por la transición CNAE-2009 $\rightarrow$ CNAE-2025 sobre series clave de
política industrial; y \emph{(d)} cruce con afiliación al Régimen General de
Seguridad Social como validación independiente de las series EPA.

El código y los datos que reproducen todos los análisis presentados están
disponibles bajo licencia MIT en
\href{https://github.com/jtamames/empleo-tic-mvp}{github.com/jtamames/empleo-tic-mvp},
con DOI Zenodo asociado. La estadística pública es objeto de filosofía
política contemporánea; este es un primer banco de pruebas empírico.

% ═══════════════════════════════════════════════════════════════════════════
\section*{Disponibilidad de datos y código}

Todo el material reproducible (datos, código de análisis, notebooks Jupyter
con el análisis estadístico formal completo, app Streamlit interactiva y este
manuscrito en LaTeX) está depositado bajo licencia MIT en el repositorio
GitHub citado, con identificador persistente en Zenodo.

\section*{Agradecimientos}

A Susana Checa (UNIE) por discusiones sobre las implicaciones jurídicas del
cambio de clasificación. Los errores remanentes son responsabilidad exclusiva
del autor.

% ═══════════════════════════════════════════════════════════════════════════
\bibliographystyle{plainnat}
\begin{thebibliography}{99}

\bibitem[Hamed and Rao(1998)]{hamed1998}
Hamed, K.~H. and Rao, A.~R. (1998).
A modified Mann-Kendall trend test for autocorrelated data.
\textit{Journal of Hydrology}, 204(1):182--196.

\bibitem[INE(2026)]{ine2026}
Instituto Nacional de Estadística (2026).
\textit{Encuesta de Población Activa, primer trimestre de 2026 --- Nota de prensa}.
INE, Madrid.

\bibitem[ONTSI(2025)]{ontsi2025}
Observatorio Nacional de Tecnología y Sociedad (2025).
\textit{Empleo Tecnológico en España. Edición 2025}.
ONTSI / Red.es, Madrid.

\bibitem[Randstad Research(2026)]{randstad2026}
Randstad Research (2026).
\textit{Mercado de trabajo en sector telecomunicaciones e IT 2025}.
Randstad España, Madrid.

\bibitem[Reglamento UE(2023)]{ue2023137}
Reglamento Delegado (UE) 2023/137 de la Comisión, de 10 de octubre de 2022,
que modifica el Reglamento (CE) n.\textsuperscript{o} 1893/2006.
\textit{Diario Oficial de la Unión Europea}, L 19/1.

\bibitem[Real Decreto 10/2025]{rd102025}
Real Decreto 10/2025, de 14 de enero, por el que se aprueba la Clasificación
Nacional de Actividades Económicas 2025 (CNAE-2025).
\textit{Boletín Oficial del Estado}, 13 (15 enero 2025), pp. 4801--5014.

\bibitem[Tamames(2024)]{tamames2024causa}
Fernández Tamames, J. (2024).
\textit{La Causa Primera: sinfonía filosófica en cinco movimientos con obertura}.
PhilArchive / Zenodo. \href{https://philarchive.org/}{philarchive.org}.

\end{thebibliography}

\end{document}
